%%%%%%%%%%%%%%%%%%%%%%%%%%%%%%%%%%%%%%%%%%%%%%%%%%%%%%%%%%%%%%%%%%%%%%%%%%%
% Englische Fassung des Uebungsblatts zu Lehreinheit 5 (Teil A und B).
% Vorlage: handouts/LE05.tex. Nur in handouts.tex mit \def\wuzsprache{en}.
%%%%%%%%%%%%%%%%%%%%%%%%%%%%%%%%%%%%%%%%%%%%%%%%%%%%%%%%%%%%%%%%%%%%%%%%%%%
\begin{uebungsblatt}{Hypotheses and Falsifiability}

\aufgabe{Falsifiable or not?}[in class, small groups, 10 min.]
For each statement, decide whether it is falsifiable \emph{in its present
form} (yes / no / only after it has been made more precise). Falsifiable
means: you can name an observation that would contradict the statement.
For every testable statement, give a concrete observation that would
refute it. Rephrase vague statements so that they become testable. (The
statements were written for this exercise.)

\begin{enumerate}
  \item[S1] \enquote{Honest people sleep better.}
  \item[S2] \enquote{Students who sleep on average at least seven hours a
    night in the week before an exam score, on average, more points in
    that exam than students who sleep less than five hours.}
  \item[S3] \enquote{Everything that happens happens for a reason, even if
    we cannot always see it.}
  \item[S4] \enquote{Astrologers can match a birth chart to the correct
    one of three personality profiles more often than chance.}
  \item[S5] \enquote{The remedy works for everyone who is open to it.
    Anyone who feels no effect was not inwardly ready.}
\end{enumerate}

\begin{center}
\renewcommand{\arraystretch}{2.4}
\begin{tabular}{|p{0.06\textwidth}|p{0.14\textwidth}|p{0.68\textwidth}|}
\hline
\textbf{No.} & \textbf{falsifiable?} & \textbf{Refutation condition or more precise version} \\
\hline
S1 & & \\ \hline
S2 & & \\ \hline
S3 & & \\ \hline
S4 & & \\ \hline
S5 & & \\ \hline
\end{tabular}
\end{center}
\end{uebungsblatt}

\begin{uebungsblattb}{2}{The Switchboard and Your Own Hypothesis}

\aufgabe{Looking back at the switchboard}[in class or homework]
\begin{enumerate}
  \item Write down the model your group built in round~1 and the number of
    trials that confirmed it.
  \item Which \emph{single} input would have put your model at greatest
    risk? Explain using Horwich's risk of refutation: a test confirms all
    the more strongly, the less likely its outcome would be without the
    model.
  \item Name two auxiliary assumptions you made without saying so while
    exploring (Duhem--Quine: what is tested is always the model
    \emph{together with} such assumptions).
\end{enumerate}
\schreiblinien{5}

\aufgabe{Null and alternative hypothesis for your research question}[start in class, due by Unit~6]
\begin{center}
\renewcommand{\arraystretch}{1.6}
\begin{tabular}{|p{0.3\textwidth}|p{0.6\textwidth}|}
\hline
My research question (from Unit~2) & \\[6mm] \hline
Type of hypothesis (difference / association / change; directional?) & \\ \hline
H\textsubscript{1} (alternative hypothesis) & \\[8mm] \hline
H\textsubscript{0} (null hypothesis) & \\[8mm] \hline
Operationalisation: concept $\to$ measurement rule & \\[8mm] \hline
Population, conditions, time frame & \\ \hline
I reject H\textsubscript{1} if \ldots & \\[8mm] \hline
Auxiliary assumptions I make & \\[6mm] \hline
\end{tabular}
\end{center}

\aufgabe{Diary of wonder}[homework]
Describe something that happened this week in which you (or someone else)
merely confirmed an assumption instead of testing it. Which trial would
really have put the assumption at risk?

\begin{lehrhinweis}[Expected answers]
\textbf{Task~1.} S1: not testable in this form, because \enquote{honest}
and \enquote{sleep better} are undefined; can be made precise, e.g.\
\enquote{Adults with higher scores on a validated honesty questionnaire
have a longer sleep duration, measured by activity tracker over two
weeks}; refuted if there is no association or a negative one. S2:
falsifiable; refuted if the group with at least seven hours scores the
same or fewer points on average. Good answers notice that such a finding
does not yet show causation (link to Unit~7). S3: not falsifiable; the
clause \enquote{even if we cannot see it} rules out every
counter-observation. S4: falsifiable; refuted if the hit rate in a
double-blind trial is not above one third. The claim was actually tested
this way \parencite{carlson1985doubleblind}, with a result at chance
level. S5: immunised; every missing effect is blamed on the person. The
claim only becomes testable if \enquote{open to it} is measured
independently and \emph{before} taking the remedy.
\textbf{Task~2.} Assessed: whether the chosen input has an outcome that
the model does \emph{not} predict with certainty, and whether auxiliary
assumptions are named, e.g.\ working lamps, full batteries, visible wiring
being the same as the actual wiring.
\textbf{Task~3.} H\textsubscript{0} and H\textsubscript{1} exclude each
other; the rejection condition names a measurement and a criterion; the
operationalisation covers every concept in the hypothesis.
\end{lehrhinweis}
\end{uebungsblattb}
